\documentclass[11pt]{article}
\usepackage[margin=1in]{geometry}
\usepackage[T1]{fontenc}
\usepackage[utf8]{inputenc}
\usepackage{lmodern}
\usepackage{microtype}
\usepackage{amsmath,amssymb,amsthm,mathtools}
\usepackage{enumitem}
\usepackage{hyperref}
\usepackage{xcolor}
\usepackage{setspace}
\usepackage{titlesec}
\usepackage{booktabs}
\usepackage{array}
\usepackage{graphicx}
\hypersetup{colorlinks=true,linkcolor=blue,citecolor=blue,urlcolor=blue}
\setstretch{1.08}
\titleformat{\section}{\large\bfseries}{\thesection.}{0.6em}{}
\titleformat{\subsection}{\normalsize\bfseries}{\thesubsection.}{0.5em}{}
\newtheorem{axiom}{Axiom}
\newtheorem{definition}{Definition}
\newtheorem{proposition}{Proposition}
\newtheorem{remark}{Remark}
\newtheorem{conjecture}{Conjecture}
\newtheorem{theorem}{Theorem}
\newcommand{\RA}{\mathrm{RA}}
\newcommand{\bdg}{\mathrm{BDG}}
\newcommand{\LLC}{\mathrm{LLC}}
\newcommand{\Pacc}{P_{\mathrm{acc}}}
\newcommand{\Pot}{\Pi}
\begin{document}
\begin{center}
{\LARGE \textbf{Relational Actualism I: Kernel, Actualization, and the Engine of Becoming}}\\[0.75em]
{\large Joshua F. Sandeman}\\[0.25em]
{\normalsize Draft updated April 13, 2026}
\end{center}

\begin{abstract}
This first paper of the three-paper Relational Actualism suite states the finite, sequential ontology of the framework and the resulting Engine of Becoming. The principal update in the present draft is the explicit adoption of \emph{Finitary Actuality}: every physically realized universe-state is represented by a finite causal graph. This axiom was already implicit in the stagewise growth dynamics, but its explicit statement clarifies the status of continuum limits, rules out completed infinities as physical states, and sharpens the treatment of severance, entropy, and cosmological nucleation developed in the companion papers. The paper presents the growing causal graph, the structured potentia, the BDG acceptance filter, and the finite-step stochastic update rule in a unified format, while distinguishing physical ontology from continuum translation.
\end{abstract}

\section{Introduction}
Relational Actualism (RA) proposes that the primitive physical object is not a timeless manifold but a growing causal graph of actualization events. The universe does not inhabit a pre-existing stage; the stage is continuously produced by the sequence of irreversible additions to the graph. Quantum potentia, classical actuality, spacetime, and macroscopic law all emerge from different aspects or limits of this single process.

The present paper is the kernel paper of the three-paper suite. Its task is narrower than a full survey: it states the ontological commitments and the sequential growth rule in the form most useful for the companion papers on matter and forces, and on gravity, cosmology, severance, and complexity. The most important revision in the present draft is that the framework now states explicitly what had already been used implicitly throughout: the universe is finite at every realized stage. This is not a stylistic preference but a physical claim.

\section{The seven axioms}
\begin{axiom}[Causal graph ontology]
At every realized stage, the universe possesses an actualized causal graph $G_n=(V_n,\prec_n)$ whose vertices are irreversible actualization events.
\end{axiom}

\begin{axiom}[Irreversibility]
Once a vertex is actualized, it is not deleted. The realized graph only grows.
\end{axiom}

\begin{axiom}[Structured potentia]
The next physically possible actualizations are not arbitrary. They are determined by the current graph through a structured possibility set $\Pot(G_n)$ of admissible single-vertex extensions.
\end{axiom}

\begin{axiom}[BDG filtering]
Candidate extensions are filtered by the four-dimensional Benincasa-Dowker score
\[
S(v,G)=1-N_1(v)+9N_2(v)-16N_3(v)+8N_4(v),
\]
and are admissible only when $S>0$.
\end{axiom}

\begin{axiom}[Stochastic actualization]
Among admissible candidates, one is selected stochastically according to the quantum measure / Born weighting appropriate to the current potentia.
\end{axiom}

\begin{axiom}[Local ledger conservation]
At every actualization vertex, the local ledger condition holds: incoming and outgoing conserved charges balance locally.
\end{axiom}

\begin{axiom}[Finitary Actuality]
Every physically realized universe-state $U_n=(G_n,\Pot(G_n))$ has finite actual graph $G_n$. Accordingly, all physically realized counts derived from $G_n$---including vertices, links, boundary sets, and admissible extensions---are finite. Infinite structures may be used as asymptotic or continuum idealizations, but are not ontologically actual.
\end{axiom}

\begin{remark}
Axiom 7 is not an add-on invented to cure a technical defect. It is the explicit statement of the discrete-first methodology already built into the finite-step growth rule. Its function is to separate ontology from approximation: completed infinities are forbidden as physical states, but large-$N$ and continuum limits remain available as mathematical descriptions of finite-but-enormous graphs.
\end{remark}

\section{Finite universe-states and the Engine of Becoming}
\begin{definition}[Universe-state]
At step $n$, the universe-state is the ordered pair
\[
U_n=(G_n,\Pot(G_n)),
\]
where $G_n$ is the finite actualized graph and $\Pot(G_n)$ is the set of admissible single-vertex extensions.
\end{definition}

\begin{definition}[Sequential growth rule]
Given finite $G_n$, a candidate past set $P\subseteq V_n$ is proposed. The new vertex $v_{n+1}$ is assigned the downward closure of $P$ as its causal past. The depth profile $(N_1,N_2,N_3,N_4)$ is computed, the BDG score is evaluated, inadmissible candidates with $S\le 0$ are rejected, and one admissible candidate is actualized stochastically to form $G_{n+1}$.
\end{definition}

\begin{proposition}[Finite potentia]
For every physically realized $U_n$, the possibility set $\Pot(G_n)$ is finite.
\end{proposition}

\begin{proof}
By Axiom 7, $G_n$ is finite. Candidate past sets are subsets of a finite set, hence finite in number. After BDG filtering, the admissible subset remains finite.
\end{proof}

\begin{remark}
This proposition matters structurally: the potentia are not an unbounded sea of possibilities but a finite, state-dependent set. The quantum measure is therefore assigned over a finite combinatorial possibility space at every realized stage, even if asymptotic continuum descriptions later approximate that space by integrals.
\end{remark}

\section{Why finitary actuality matters}
The explicit statement of finitary actuality reorganizes several recurrent themes in the RA programme.

\subsection{Completed infinities are not physical states}
A continuum limit may be useful as an approximation, but it is not ontologically basic. There is no realized state with infinitely many vertices, infinitely many links, or infinitely many physically instantiated degrees of freedom. This has immediate consequences:
\begin{enumerate}[label=(\alph*)]
    \item ultraviolet divergences are not physical because the graph is discrete and finite at every stage;
    \item singularities are not physical because graph growth terminates or redirects through severance rather than producing infinite curvature;
    \item vacuum-catastrophe style infinities are not physical because unactualized potentia do not count as realized metric sources;
    \item infinite-time pathologies such as Boltzmann-brain arguments lose their direct ontological force because infinite future duration is not a realized graph state but an asymptotic idealization.
\end{enumerate}

\subsection{Physics versus translation}
The distinction between discrete law and continuum translation is now explicit.
\begin{itemize}
    \item Discrete physics asks what the graph is doing: which vertices exist, which links are present, which counts are finite, and what local rules govern growth.
    \item Translation asks how those finite counts appear when described in metric language: proper time, curvature, horizon area, entropy densities, and so forth.
\end{itemize}
This distinction will be crucial in Paper III, where the severance entropy is defined as a finite count of severed links and only subsequently compared with Bekenstein-Hawking entropy.

\section{The BDG filter and acceptance kernel}
In the Poisson-CSG approximation, the BDG acceptance probability at density $\mu$ is
\[
\Pacc(\mu)=\mathbb{P}(S>0), \qquad S=1-N_1+9N_2-16N_3+8N_4,
\]
where each $N_k$ is Poisson distributed with parameter $\lambda_k=\mu^k/k!$. The resulting acceptance kernel captures how the finite local graph is distorted by the BDG filter.

The important conceptual point for the present paper is not the detailed shape of $\Pacc(\mu)$ but its status: it is a probability assigned to finite local profiles in a finite-step graph growth process. Even when one studies its asymptotic behavior or large-$N$ limit, the underlying ontology remains finitary.

\section{Severance preview}
A causal severance is a partition of the realized graph into subgraphs whose cross-links are cut or blocked. In Paper III, the entropy of such a severance will be defined as the number of irreducible cross-severance links removed by the severance. The introduction of Axiom 7 already prepares that analysis. Since every realized $G_n$ is finite, every physically realized severance entropy is automatically finite. The remaining work is not to rescue the observable from infinities but to identify its invariant content and then study its continuum translation.

\section{Conclusion}
The explicit adoption of Finitary Actuality clarifies the deepest methodological commitment of Relational Actualism: the discrete graph is reality, while continuum structures are descriptive idealizations of finite graphs. This axiom strengthens the Engine of Becoming rather than altering it. It renders the universe-state finite at every realized stage, sharpens the distinction between ontology and translation, and provides the correct foundation for the severance, entropy, and cosmological arguments of the companion papers.

\section*{References}
\begin{enumerate}[leftmargin=1.5em]
\item J. F. Sandeman, \emph{Relational Actualism: The Engine of Becoming}, draft suite notes, 2026.
\item D. M. T. Benincasa and F. Dowker, ``The scalar curvature of a causal set,'' \emph{Phys. Rev. Lett.} \textbf{104}, 181301 (2010).
\item R. D. Sorkin and collaborators on classical sequential growth and causal set dynamics.
\end{enumerate}
\end{document}
